\documentclass[dvipdfmx,autodetect-engine,a4paper,10pt]{jsarticle}

\usepackage[top=22mm,bottom=22mm,left=22mm,right=22mm]{geometry}
\usepackage{amsmath,amssymb,bm}
\usepackage{multicol}
\usepackage{enumitem}
\usepackage{titlesec}
\usepackage[dvipdfmx]{hyperref}

\pagestyle{plain}
\setlength{\parindent}{1em}
\setlength{\parskip}{0pt}
\setlength{\columnsep}{8mm}
\renewcommand{\baselinestretch}{0.95}

\titleformat{\section}
  {\normalfont\bfseries\normalsize}
  {\thesection.}{0.7em}{}
\titlespacing*{\section}{0pt}{0.7\baselineskip}{0.2\baselineskip}

\newcommand{\E}{\mathbb E}
\newcommand{\Prb}{\mathbb P}
\newcommand{\indep}{\mathrel{\perp\!\!\!\perp}}

\newcommand{\papertitle}{%
  ベクトル値欠測変数に対する欠測IV識別
}

\newcommand{\authorblock}{%
  \begin{tabular}{lll}
    慶應義塾大学 & 本田 昌宏 & HONDA Masahiro\\
    慶應義塾大学 & 星野 崇宏 & HOSHINO Takahiro\\
  \end{tabular}
}

\begin{document}

\vspace*{12mm}

\begin{center}
  {\Large\bfseries \papertitle\par}
  \vspace{6mm}
  {\normalsize \authorblock}
\end{center}

\vspace{4mm}

\begin{multicols}{2}

\section{はじめに}
ここに導入を書く。問題設定，既存研究との差分，本稿の貢献を短く述べる。

\section{識別設定}
ここに記号と仮定を書く。たとえば，欠測対象を
\(\bm Y=(Y_1,\ldots,Y_K)'\)，観測パターンを
\(\bm D=(D_1,\ldots,D_K)'\)，完全観測インジケータを
\[
R=\mathbf 1\{\bm D=\bm 1\}
\]
とする。

\section{主結果}
ここに主定理を書く。中心となる条件は
\[
R\indep \bm Z\mid \bm Y
\]
と vector completeness であり，選択確率
\[
\pi(\bm y)=\Prb(R=1\mid \bm Y=\bm y)
\]
は
\[
\E\left\{\frac{R}{\pi(\bm Y)}\mid \bm Z\right\}=1
\]
を通じて識別される。

\section{まとめ}
ここに結論を書く。完全ケースの選択確率を補正することで，
母集団の \(P(\bm Y,\bm Z)\) および \(P(\bm Y)\) を回復できることを述べる。

\section*{参考文献}
\begin{enumerate}[label={[\arabic*]},leftmargin=1.7em,itemsep=0pt,topsep=0pt]
  \item X. d'Haultfoeuille, A new instrumental method for dealing with
  endogenous selection, \textit{Journal of Econometrics} 154, 1--15, 2010.
\end{enumerate}

\end{multicols}

\end{document}
